\label{sec:conclusiones}

El desarrollo de esta experiencia permitió cumplir satisfactoriamente con los objetivos planteados, centrados en la construcción, análisis y validación de un amperímetro analógico multirrango, a pesar de las adaptaciones metodológicas forzadas por fallas en el equipamiento (potenciómetros). 

El cambio de metodología, utilizando una resistencia limitadora fija y variando la tensión de la fuente, demostró que es posible contrastar y verificar la linealidad del instrumento. Sin embargo, al depender de una única resistencia de calibración ($R_c$) optimizada empíricamente para el rango de 100 mA, se evidenció que los rangos inferiores sufren un error sistemático inherente debido a que no operan en su punto óptimo de sintonía.

En cuanto a la linealidad, se comprobó experimentalmente que el diseño basado en divisores resistivos (shunt) mantiene un comportamiento lineal casi perfecto ($R^2 > 0.996$ en todos los rangos). No obstante, el fuerte "offset" o intercepto observado en las curvas de calibración de los rangos superiores confirma la presencia de un marcado error de exactitud. Esto subraya que la linealidad (precisión y proporcionalidad) no garantiza la exactitud si la calibración del punto cero u offset es inadecuada.

El análisis del circuito de carga en paralelo fue fundamental para ilustrar el \textit{efecto de carga} y la naturaleza intrusiva de la medición de corriente. Al utilizar la fuente de 10 V con las nuevas resistencias ($1\text{ k}\Omega, 330\text{ }\Omega, 220\text{ }\Omega$), se observó una perturbación extrema del circuito. Dado que el amperímetro presentaba una alta resistencia interna (producto de las adaptaciones del circuito de calibración), la mera inserción del instrumento duplicó la resistencia total en varias ramas, provocando caídas de la corriente real de hasta un 60\% respecto al valor teórico esperado (ver Tabla \ref{tab:paralelo_error_carga}). Esto resalta drásticamente la importancia crítica de diseñar amperímetros con una impedancia interna que tienda a cero, para evitar desvirtuar la magnitud que se intenta medir.

Finalmente, el análisis de error y su posterior mitigación matemática validaron la robustez del modelo lineal. Aunque el amperímetro presentaba errores brutos de lectura superiores al 36\% debido a su offset y efecto de carga, el uso de la ecuación inversa de la recta de calibración ($I_{a,\text{corr}} = \frac{I_a - b}{m}$) permitió reconstruir el valor verdadero de la corriente con un error corregido inferior al 2\% en todos los casos. Esto demuestra que, siempre y cuando un instrumento sea rigurosamente lineal y repetible, sus deficiencias de exactitud sistemática pueden ser suprimidas exitosamente mediante una caracterización y compensación matemática adecuada.